\scanpage{245}
\mspage{125}

On voit tout de suite

(76) La donnée du graphe à opérateurs homogène $(S,A_0\underset{\deleted{\unsure{\gg|=}}}{\subset}\mathfrak P_{\{1,2\}}(S), H\subset\unsure{\mathfrak S|\mathcal G}_S)$, muni de l'arête orientée $(s_0,t_0)$, équivaut à la donnée du groupe $H$, des deux sous-groupes $U_0,U_1$ de $H$, et de $\sigma\in H$ satisfaisant les relations (75), et de plus

\[
\bigcap_{u\in H}\mathrm{int}(u)U_0=\{1\}
\tag*{(77)}
\]

qui exprime que l'opération de $H$ sur l'espace homogène $H/U_0$ est fidèle.

On récupère la situation géométrique, \unsure{et|de|ou} la situation algébrique, par

\[
\begin{aligned}
&S\simeq H/U_0,\quad s_0=U_0,\quad t_0=\sigma\mathbin{.}H_0\quad\text{(translaté à gauche)}\\
&\overrightarrow{A_0}\simeq H/\unsure{\mathfrak z|\zeta}
\end{aligned}
\tag*{(78)}
\]

\deleted{Dans} L'énoncé (76) n'est correct que si, ou bien on sous-entend que parmi les données géométriques il y a aussi le \uline{choix} de $\sigma$ satisfaisant (74), soit que, dans la traduction algébrique, $\sigma$ \deleted{n'est donné que} \inserted{(qui normalise $\mathfrak z_H$, nécessairement)} n'est donné que \textquotedblleft{}modulo $\mathfrak z_H$\textquotedblright{} (i.e. que ce qui est donné est la classe \inserted{$\sigma\mathfrak z_H=\mathfrak z_H\sigma$} \deleted{de} $\sigma$ dans $\mathrm{Norm}_H(\mathfrak z_H)/\mathfrak z_H$). Cette ambiguïté \unsure{disparait|disparaît} quand on suppose que $H$ est \uline{simplement} transitif sur l'ens \inserted{$\overrightarrow{A_0}$} des arêtes orientées ; \deleted{ce que}

\[
\begin{aligned}
&H\text{ simplement transitif sur l'ens }\overrightarrow{A_0}\\
&\Longleftrightarrow\mathfrak z_H(\overset{\mathrm{def}}{=}U_0\cap U_1)=\{1\},\quad\deleted{\text{ce q}}
\end{aligned}
\tag*{(79)}
\]

Condition qui implique l'unicité de $\sigma$, et

\[
\sigma^2=1 .
\tag*{(80)}
\]

% archive sheet number: 244
